\chapter{Minería y Correlación de Bitácoras del Sistema}

Las bitácoras de auditoría del sistema operativo (\texttt{/var/log/messages}, \texttt{/var/log/auth.log}, \texttt{/var/log/secure} y \texttt{/var/log/xferlog}) constituyen el registro documental más elocuente de la actividad de una máquina Unix. Mientras que el análisis del sistema de archivos revela los cambios materiales que persistieron en el disco, las bitácoras permiten reconstruir la interacción dinámica entre el atacante y los servicios de red: las conexiones entrantes, las solicitudes de autenticación fallidas y exitosas, y las transferencias de archivos.

En este capítulo se realiza la minería forense de los registros del servidor \texttt{victima1}, desentrañando el vector de compromiso inicial a través de cuentas de usuario vulneradas y las maniobras de transferencia de código malicioso mediante el protocolo FTP.

\section{Análisis de Bitácoras de Autenticación (\texttt{auth.log} / \texttt{secure})}

En las distribuciones GNU/Linux de la época, las solicitudes de inicio de sesión a través de consola, terminales remotas (\textit{telnet}, \textit{rlogin}) y el protocolo SSH eran procesadas mediante el subsistema \textbf{PAM} (\textit{Pluggable Authentication Modules}) y registradas en \texttt{/var/log/auth.log}.

La inspección cronológica de los registros del 8 de abril de 2004 ---el mismo día en que la línea de tiempo reveló la implantación del \textit{rootkit}--- arrojó la secuencia reveladora documentada en el Listado~\ref{lst:authlog_intrusion}.

\begin{lstlisting}[style=terminal,caption={Registro de autenticación auth.log durante la intrusión inicial del 8 de abril.},label={lst:authlog_intrusion}]
Apr  8 05:29:22 victima1 sshd[1232]: Did not receive identification string from 198.18.16.225
Apr  8 05:30:20 victima1 PAM_unix[1234]: check pass; user unknown
Apr  8 05:30:20 victima1 PAM_unix[1234]: authentication failure; (uid=0) -> **unknown** for login service
Apr  8 05:30:22 victima1 login[1234]: FAILED LOGIN (1) on pts/2 from dsl-198-18-16-225.isp-remoto.net FOR UNKNOWN
Apr  8 05:30:29 victima1 PAM_unix[1234]: (login) session opened for user proa by (uid=0)
Apr  8 05:31:37 victima1 newgrp[1242]: unknown gid: 523
Apr  8 05:31:56 victima1 passwd[1252]: can't change pwd for `root'
Apr  8 05:33:02 victima1 PAM_unix[1262]: (login) session opened for user proa by (uid=0)
Apr  8 05:34:58 victima1 passwd[4388]: can't change pwd for `root'
Apr  8 05:36:23 victima1 PAM_unix[4402]: (login) session opened for user proa by (uid=0)
Apr  8 05:36:49 victima1 PAM_unix[4404]: authentication failure; proa(uid=523) -> irc for su service
Apr  8 05:36:51 victima1 su[4404]: pam_authenticate: Authentication failure
Apr  8 05:36:51 victima1 su[4404]: - pts/5 proa-irc
Apr  8 05:43:05 victima1 PAM_unix[4588]: (login) session opened for user proa by (uid=0)
\end{lstlisting}

\subsection{Disección Forense de la Intrusión Inicial}

El análisis de esta secuencia temporal suministró evidencia pericial cardinal sobre el modus operandi del intruso:
\begin{enumerate}
    \item \textbf{05:29:22~h (Sondeo de puertos):} El atacante, conectado desde la dirección IP externa \texttt{198.18.16.225} (\texttt{dsl-198-18-16-225.isp-remoto.net}), ejecutó un escaneo inicial contra el servicio SSH (\texttt{sshd[1232]}).
    \item \textbf{05:30:22~h (Fuerza bruta y acceso inicial):} Tras un intento fallido de autenticación con un usuario inexistente (\texttt{UNKNOWN}), apenas siete segundos después (a las 05:30:29~h) logró abrir exitosamente una sesión interactiva autenticándose como el usuario legítimo \texttt{proa} (UID 523).
    \item \textbf{05:31:56~h (Intento prematuro de cambiar la contraseña de root):} Menos de dos minutos después de ingresar, el intruso ejecutó el comando \texttt{passwd root} en un intento ingenuo por cambiar la clave del superusuario. Al carecer de privilegios administrativos en ese instante, el sistema rechazó la solicitud (\texttt{can't change pwd for `root'}).
    \item \textbf{05:36:49~h (Sondeo de cuentas de servicio con su):} El atacante intentó transitar mediante el comando \texttt{su} hacia una cuenta de servicio denominada \texttt{irc}, buscando probablemente acceder a procesos de red preexistentes.
\end{enumerate}

Esta conducta demuestra que el atacante disponía de credenciales comprometidas de antemano o que obtuvo acceso mediante un ataque de fuerza bruta por diccionario sobre contraseñas débiles, explorando agresivamente los permisos del sistema operativo.

\section{Auditoría de Transferencias: El Registro del Servidor FTP (\texttt{xferlog})}

Una vez afianzada una cuenta de acceso, la prioridad operativa del atacante consistió en introducir su arsenal de herramientas. Para ello empleó el servicio de transferencia de archivos (FTP), cuyas actividades quedaron registradas en el archivo canónico \texttt{/var/log/xferlog}.

El Listado~\ref{lst:xferlog_analysis} expone las transferencias ejecutadas hacia el servidor.

\begin{lstlisting}[style=terminal,caption={Registros de transferencia de paquetes de malware recuperados de xferlog.},label={lst:xferlog_analysis}]
Tue Apr 08 05:45:10 2004 1 198.18.16.225 152841 /home/proa/b.tar.gz b _ i r proa ftp 0 * c
Tue Apr 08 05:47:22 2004 1 198.18.16.225 845201 /home/proa/psy.tar.gz b _ i r proa ftp 0 * c
Tue Apr 08 05:52:14 2004 1 198.18.16.225  45210 /home/proa/energy.tar.gz b _ i r proa ftp 0 * c
Tue Apr 08 06:14:02 2004 1 198.18.16.225 612400 /home/proa/rk.tar.gz b _ i r proa ftp 0 * c
\end{lstlisting}

\subsection{Anatomía y Estructura del Registro xferlog}

Cada línea del archivo \texttt{xferlog} se ajusta a la especificación estándar del demonio FTP (\textit{wu-ftpd} / \textit{vsftpd}), desglosada en el Cuadro~\ref{tab:xferlog_fields}.

\begin{table}[htbp]
\centering
\small
\begin{tabular}{lll}
\toprule
\textbf{Campo} & \textbf{Valor en Listado~\ref{lst:xferlog_analysis}} & \textbf{Significado Forense} \\
\midrule
Fecha y Hora & \texttt{Tue Apr 08 05:45:10} & Momento de finalización de la transferencia. \\
Tiempo de transferencia & \texttt{1} & Duración en segundos de la subida. \\
Host remoto & \texttt{198.18.16.225} & Dirección IP externa del atacante. \\
Tamaño del archivo & \texttt{152841} & 152.8~KB transferidos. \\
Nombre del archivo & \texttt{/home/proa/b.tar.gz} & Archivo comprimido con código fuente malicioso. \\
Tipo de transferencia & \texttt{b} & Modo binario (\textit{binary}). \\
Dirección de flujo & \texttt{i} & Flujo entrante hacia el servidor (\textit{incoming / upload}). \\
Modo de acceso & \texttt{r} & Usuario autenticado local real (\textit{real user}). \\
Nombre de usuario & \texttt{proa} & Cuenta empleada para autenticar en el FTP. \\
Servicio & \texttt{ftp} & Demonio receptor. \\
Estado & \texttt{c} & Transferencia completada exitosamente (\textit{complete}). \\
\bottomrule
\end{tabular}
\caption{Desglose de campos del formato estándar xferlog.}
\label{tab:xferlog_fields}
\end{table}

\subsection{Correlación de los Paquetes Descargados}

El cotejo de los nombres y tamaños de archivo con la evidencia recolectada en los capítulos previos completó el mapa del incidente:
\begin{itemize}
    \item \texttt{b.tar.gz}: Archivo comprimido que contenía los scripts de instalación y los ficheros de configuración encubiertos (\texttt{.libn}, \texttt{.libpt}).
    \item \texttt{psy.tar.gz}: Código fuente del servidor de retransmisión \textit{psyBNC}, desempaquetado y compilado posteriormente mediante \texttt{gcc}.
    \item \texttt{energy.tar.gz}: Distribución del bot IRC \textit{Energyb}.
    \item \texttt{rk.tar.gz}: Paquete del \textit{Linux Rootkit} que contenía los binarios modificados de \texttt{ps}, \texttt{ls} y \texttt{netstat}.
\end{itemize}

\section{Intentos de Evasión y Desfases Horarios}

Al cruzar los registros de \texttt{auth.log} con \texttt{xferlog} y las marcas temporales del sistema de archivos, el equipo pericial detectó intentos deliberados de manipulación de bitácoras:
\begin{itemize}
    \item \textbf{Truncamiento de bitácoras web y de correo:} El atacante intentó en varias ocasiones vaciar los registros de acceso mediante redirecciones nulas (\texttt{> /var/log/xferlog}), aunque la activación preventiva del atributo de solo adición (\texttt{chattr +a}) sobre los archivos de log bloqueó la supresión de las trazas anteriores.
    \item \textbf{Desfase horario (\textit{clock skew}):} Se identificaron pequeñas discrepancias de varios minutos entre las marcas de tiempo locales del servidor y los sellos de red reportados por clientes remotos, una contingencia típica en servidores de 2004 que no contaban con sincronización horaria continua mediante el protocolo NTP (\textit{Network Time Protocol}).
\end{itemize}

\section{Conclusiones de la Minería de Registros}

La investigación de bitácoras consolidó la reconstrucción pericial de la intrusión externa:
\begin{itemize}
    \item El punto de entrada inicial fue una brecha de autenticación en la cuenta \texttt{proa} a las 05:30~h del 8 de abril desde una dirección IP residencial externa.
    \item Entre las 05:45~h y las 06:14~h, el intruso cargó vía FTP los cuatro paquetes de software malicioso que conformaron la infraestructura del ataque.
    \item Menos de diez horas después (a las 16:02~h, como demostró la línea de tiempo del Capítulo 8), los paquetes fueron desempaquetados e instalados, consumando el compromiso del superusuario.
\end{itemize}
